\chapter{The chain}

\textbf{Purpose:} Give the wall an exact law and decompose it, so a depth figure is computed rather than searched for. \textbf{Scope:} \texttt{src/\allowbreak{}bench/\allowbreak{}bench\_\allowbreak{}words.\allowbreak{}cu}, \texttt{src/\allowbreak{}bench/\allowbreak{}bench\_\allowbreak{}depth\_\allowbreak{}cuda.\allowbreak{}cu}.

\section{The wall is the chain length}

A one-bit difference in the chaining state stops being visible to a per-word Hamming statistic at a definite depth. That depth is not a property of the mixing functions. It is the number of words on the shift chain.

The round function generalises to any even word count: majority and \(\Sigma_0\) read the low half's head, choice and \(\Sigma_1\) read the high half's head, and the two halves join the way the standard joins them. At eight words this is SHA-256 exactly rather than a model of it, which makes that row the control.

\begin{center}
\begin{tabular}{rrr}
\toprule
words & wall & \(\text{words} + 2\) \\
\midrule
8 & 10 & 10 \\
10 & 12 & 12 \\
12 & 14 & 14 \\
14 & 16 & 16 \\
16 & 18 & 18 \\
\bottomrule
\end{tabular}
\end{center}

Eight words reproduces the wall of \(10\) measured independently by the depth bench, and the other rows are read against that row.

\textbf{The offset is structure, not sample size.} The same table returns identical at \(2^{20}\), \(2^{24}\) and \(2^{28}\) pairs, a two-hundred-and-fifty-six-fold range with no drift. A cliff predicts that and a slope does not: a collapse completing inside one round cannot be pushed by adding samples.

\section{Why plus two}

Only two slots of the chain are mixed at all. Slot zero takes \(T_1 + T_2\) and slot \(\mathrm{half}\) takes \(d + T_1\). Every other slot is \(\mathrm{state}[i] = \mathrm{state}[i-1]\), a copy that transports a difference without touching it. So the tail of the chain holds whatever the join produced \(\mathrm{count} - 1\) rounds earlier, and the wall should be a mixing depth plus a transport delay rather than one quantity.

Reading each slot's own death depth instead of the largest across slots shows it. Every slot from the \(e\) join to the tail dies at exactly \(k + 3\):

\begin{center}
\begin{tabular}{rll}
\toprule
words & slots from the \(e\) join & dies at \\
\midrule
8 & 4, 5, 6, 7 & 7, 8, 9, 10 \\
10 & 5, 6, 7, 8, 9 & 8, 9, 10, 11, 12 \\
12 & 6, 7, 8, 9, 10, 11 & 9, 10, 11, 12, 13, 14 \\
\bottomrule
\end{tabular}
\end{center}

The last slot is \(\mathrm{count} - 1\), so \(\mathrm{wall} = (\mathrm{count} - 1) + 3 = \mathrm{count} + 2\). The \(3\) is the round function's mixing depth, how long \(T_1\) and \(T_2\) take to saturate a freshly computed word, and it does not depend on the word count. That is exactly why the wall tracks the count one for one. The \(-1\) is that the chain is \(\mathrm{count} - 1\) hops from the join to the tail rather than \(\mathrm{count}\).

\section{The two halves are different chains}

The \(a\) half was read as irregular before. It is not irregular. Both halves are chains of copies below one mixed head, and both follow an exact law; the laws differ in how they depend on the chain length, and comparing a count-dependent offset against a count-independent expectation makes one of them look like noise.

\begin{center}
\begin{tabular}{rlll}
\toprule
words & \(a\) half & \(e\) half & apart \\
\midrule
8 & slot \(+\, 6.00\) & slot \(+\, 3.00\) & 3.00 \\
10 & slot \(+\, 7.00\) & slot \(+\, 3.00\) & 4.00 \\
12 & slot \(+\, 8.00\) & slot \(+\, 3.00\) & 5.00 \\
14 & slot \(+\, 9.00\) & slot \(+\, 3.00\) & 6.00 \\
16 & slot \(+\, 9.75\) & slot \(+\, 3.00\) & 6.75 \\
\bottomrule
\end{tabular}
\end{center}

\textbf{The \(e\) half is slot \(+\,3\) at every count, exactly. The \(a\) half is slot \(+\, \mathrm{half} + 2\)}, giving six, seven, eight and nine for halves of four, five, six and seven. The \(9.75\) at sixteen words is one slot misreading where the law gives ten.

So the \(e\) half's saturation depth does not depend on the chain length and the \(a\) half's does, growing one round per added word pair.

\textbf{Which half sets the wall.} The \(a\) half's tail is slot \(\mathrm{half}-1\) and dies at \((\mathrm{half}-1) + \mathrm{half} + 2 = \mathrm{count} + 1\). The \(e\) half's tail is slot \(\mathrm{count}-1\) and dies at \(\mathrm{count} + 2\). The \(e\) half tail is one round deeper, so the wall comes from the \(e\) half even though the \(e\) half saturates faster per slot. It wins by extending to the end of the chain rather than by being slower.

The candidate mechanism is unconfirmed and is recorded as a candidate: \(a\) takes \(T_1\) plus \(T_2\), both halves fully mixed, while \(e\) takes \(d\), a raw copy of the \(a\) half's tail, plus \(T_1\). Twice the mixing into \(a\) is consistent with it saturating sooner per slot, and \(e\) reading a raw copy is consistent with its offset not tracking the chain. Neither is measured; only the laws are.

\section{Six words is degenerate}

Six words does not wall within thirty rounds. At three words per half, choice reads the entire high half and majority the entire low half, so neither function selects among words any more. It is excluded rather than explained away.
